\documentclass[11pt]{article}
\usepackage[margin=1in]{geometry}
\usepackage{amsmath,amssymb}
\usepackage{array,booktabs}
\usepackage{tikz}
\usepackage{enumitem}

\setlength{\parindent}{0pt}
\setlength{\parskip}{6pt}
\begin{document}

\section*{Problem 15}

\textbf{Question.} At the start of an Ebola epidemic, $0.2\%$ of the population had the virus. If a random sample of 1000 residents is taken, find the chances of finding exactly 5 infected people and at most 3 infected people.

\textbf{Solution.} Let
\[
X=\text{number of Ebola-infected people in the sample}.
\]
Since $0.2\%$ of the population had the virus,
\[
p=0.002, \qquad n=1000.
\]
Thus,
\[
X\sim \operatorname{Binomial}(1000,0.002).
\]
Since $n$ is large and $p$ is small, use a Poisson approximation with
\[
\lambda=np=1000(0.002)=2.
\]
Therefore,
\[
X\approx \operatorname{Poisson}(2).
\]

\textbf{Exactly 5 infected people:}
\[
P(X=5)=\frac{e^{-2}2^5}{5!}=\frac{e^{-2}(32)}{120}\approx 0.0361.
\]
Therefore,
\[
\boxed{P(X=5)\approx 0.0361}.
\]

\textbf{At most 3 infected people:}
\[
P(X\le 3)=P(X=0)+P(X=1)+P(X=2)+P(X=3).
\]
Using the Poisson pmf,
\[
P(X\le 3)=e^{-2}\left(\frac{2^0}{0!}+\frac{2^1}{1!}+\frac{2^2}{2!}+\frac{2^3}{3!}\right).
\]
Thus,
\[
P(X\le 3)=e^{-2}\left(1+2+2+\frac{8}{6}\right)=e^{-2}\left(\frac{19}{3}\right)\approx 0.8571.
\]
Therefore,
\[
\boxed{P(X\le 3)\approx 0.8571}.
\]

\newpage

\end{document}
